\section{Rによる構造方程式モデリング}
\subsection{授業内容}
\subsubsection{科目の中でのこのコマの位置づけ}
これまでの流れと同じで，統計技術の理論を知っただけではなく，自分で実際に計算できる演習を経てこそ理解が深まるということから，本講ではRをつかって実際に構造方程式モデリングを解くことを演習的に学ぶ。構造方程式モデリングを実装するパッケージは複数あるが，最も応用範囲がひろい\texttt{lavaan}パッケージを用いることにする。

まずは観測変数だけからなる簡単なパス解析を行う。データの入力の仕方，方程式の設定，関数の使い方などを一通り習得する。
続いて潜在変数を含んだモデルによる解析を行う。モデルの適合度や修正指数を参考に，徐々にモデルを書き換えていく手順を学ぶ。注意すべきは，適合度を上げることが目的になって，不自然な仮定やパスをおいてしまうことである。あくまでも具体的かつ妥当なモデリングを心がけるべきである。

オプショナルな設定になるが，観測変数がカテゴリカルである場合や，推定方法の選択なども確認する。
最後に，R以外の統計パッケージによる構造方程式モデリングの実践例がいくつか紹介される。

\subsubsection{コマ主題細目}
\begin{description}
	\item[方程式の入力] まずは観測変数同士の関係をパスでつなぐモデルで練習する。パス解析は回帰分析の繰り返しで実行することもできるが，構造方程式モデリングによってパスの繋がりを1つのモデルで表現し，適合度も統一できるなどの利点がある。観測変数だけからなるモデルの結果と，実際に\texttt{lm}関数で実行した結果と比較すると良い。またパッケージにもよるが，自動的にパスダイアグラムを描画してくれるものもある。方程式とそのパスダイアグラムによる表現の対応を確認する。

	      \begin{flushright}
		      $\to$\textcite{Kosugi2014Mplus}のPp.55--60
	      \end{flushright}

	\item[測定モデルの実践] 因子分析モデルをSEM上で実行してみる。探索的因子分析と違い，どの項目にどの因子が影響しているかを固定したモデリングが可能であり，この検証的因子分析による結果と，いわゆる因子分析関数との結果を比較することで，パスが引かれていないところはその係数が$0$であるという強い仮定をおいていることを確認する。また尺度作成の観点からは，検証的因子分析をすることで因子的妥当性や弁別的妥当性，収束的妥当性を検討することもできる。さらに同じモデルを別のデータに適用することで多母集団同時分析を行うことになる。このように，モデルの暗黙の過程や，モデルとデータの適合という側面にとくに注意する。さらに測定モデルと測定モデルをつなぐ，構造方程式を扱ったモデルへと拡張する。

	      \begin{flushright}
		      $\to$\textcite{Kosugi2014Mplus}のPp.87--90
	      \end{flushright}

	\item[実践上の注意点] ここまでを通じて，一通りモデルを作成できるようになった。とくに構造方程式を踏まえると，心理的実体同士の関係を描画したと解釈できるため，心理学的概念間の関係を記述できることは魅力的に映るかもしれない。しかしデータを越えての解釈はご法度であり，潜在変数が心理的実在であるかどうかの議論は，理論的背景や測定の適切さ，標準化されないスコアが実際にどのように変化すれば何が言えるのか，といったところに一足飛びに行かぬよう注意する必要がある。またモデル改良のステップにおいて，適合度や修正指数を過度に参照していないか，注意する必要がある。

	\item[そのほかの統計パッケージ] 構造方程式モデリングの利点は，モデルを可視化したことにもある。たとえばAMOSはGUIでモデルを作成できる。他にMplusはカテゴリカルな変数にも対応しているし，高度に複雑なモデルであっても表現が可能である。
\end{description}
\subsubsection{キーワード}
\begin{itemize}
	\item 測定方程式
	\item 構造方程式
	\item 潜在変数を含んだモデル
	\item 多母集団同時分析
	\item 適合度
\end{itemize}
\subsection{授業情報}
\paragraph{コマの展開方法}
講義
\subsubsection{予習・復習課題}
\paragraph{予習} 構造方程式モデルは，数式レベルでの理解は難しいが実際は統合的なものであり，回帰分析や因子分析をその下位モデルとして含んでいる。改めて，回帰分析や因子分析を単体で行った場合にどのような出力がなされるのか確認しておくと，同じものを構造方程式で実践したときの違いが明確に意識できるようになる。
\paragraph{復習} さまざまなモデルを試すなかでは，エラーや警告が出ることもある。そうしたエラーや警告の意味を理解し，またそれに対応するためにはどのような方法が取れるかを考える必要がある。まずは手元のデータを用いて，これらの練習を行うと良い。